\subsection{A finite-bit oracle for the carrier family}
\label{sec:pn-family-oracle}

The direct implementation in Theorem~\ref{thm:pn-family-query-law} has
a factor $D+1$ in its bit work. This subsection removes that factor by
preprocessing the two-mean function, which gives the bit bound in
Theorem~\ref{thm:main-carrier}.

\begin{theorem}[The carrier family's bit cost]
\label{thm:pn-family-bit-law}
Under the hypotheses of Theorem~\ref{thm:pn-family-query-law},
weight-only preprocessing of time and storage
$O_s((D+1)n^2B^6)$ gives an exact sampler with expected online
bit work
\[
 O_s\!\left(\min\{r_D^{2\nu_\eta},n\}B^4\right).
\]
The lower query bound of that theorem therefore matches this
bit bound up to the displayed factor.
\end{theorem}

\paragraph{Proof idea.}
The normalized recurrence has a complex neighborhood whose radius
shrinks as $r_D^{-\nu_\eta}$. Preprocessing covers the two unknown
means with Taylor patches at that scale. A patch contains only
polynomially many precision bits, since the recurrence propagates
rounding error through factors $1+O_s(q_j)$. Cached values answer
almost every factory comparison. An unusually close comparison
recomputes the recurrence; its probability pays for the depth.

\begin{proof}
Write $r=r_D$, $\nu=\nu_\eta$, and retain $h,C_0,C_u,M$ from
the preceding proof. We first establish a uniform complex
neighborhood. Fix a real parameter pair $(z_0,w_0)$ and its real
trajectory. Perturb each parameter by at most $\delta$ in the
complex plane, and let $e_j$ be the absolute perturbation in
$u_j$. The local derivative bounds and Taylor's formula give,
as long as $e_j+2\delta\le h/2$,
\[
 e_{j+1}\le L_j e_j+C_0q_j\delta
                  +C_2q_j(e_j+2\delta)^2,
 \qquad L_j=1+\nu q_j+C_uq_j/r_j,
\]
where a fixed $C_2=100C_0$ suffices. These bounds follow either
by integrating the second derivative along the complex segment
or by its convergent local Taylor series. Enlarge $C_0$ by two
if needed to cover the sum of the direct parameter derivatives.

Set $K=\max\{2,2C_0/\nu\}$ and $d_j=e_j+K\delta$.
Since $L_j-1\ge\nu q_j$,
\[
 d_{j+1}\le L_jd_j+C_2q_jd_j^2.
\]
Let $P_j=\prod_{i<j}L_i$. The preceding product estimate gives
$P_j\le Mr_j^\nu$. For $v_j=d_j/P_j$ this implies
\[
 v_{j+1}\le v_j+C_2q_jP_jv_j^2,
 \qquad
 \sum_{j<D}q_jP_j\le Mr^\nu/\nu.
\]
The last inequality follows from
$\sum q_jr_j^\nu\le(r^\nu-1)/\nu$.
The elementary bootstrap for $v_{j+1}\le v_j+a_jv_j^2$
says $v_j\le2v_0$ whenever $4v_0\sum a_j\le1$:
under this bound the total increment is at most
$4v_0^2\sum a_j\le v_0$. Taking a strict inequality first
and then continuity handles the boundary.

Choose a positive rational constant $c_s$ below
\[
 \min\left\{\frac{\nu}{16C_2M(K+1)},
              \frac{h}{16M(K+1)},\frac h{16}\right\},
 \qquad \delta=c_s r^{-\nu}.
\]
Uniform choices in $\eta\in[1/4,1/2]$ are possible because
$\nu\ge\sqrt2s/c-1>0$. Then
$e_j\le2M(K+1)\delta r_j^\nu\le h/8$.
The local analytic extensions remain valid at every step.
This bootstrap, first on a slightly smaller polydisk and then
by continuation, proves that every intermediate state and the
head are holomorphic on the parameter polydisk of radius
$\delta$. Their moduli there are at most two.

Choose a dyadic mesh spacing $a$ with
$\delta/32<a\le\delta/16$. Cover $[-1,1]^2$ by grid squares
with a real corner as center; truncate boundary squares at one.
There are $O_s(r^{2\nu})$ squares. At a center $(z_*,w_*)$,
use scaled coordinates
$z=z_*+2aU$, $w=w_*+2aV$. A point in its square has
$|U|,|V|\le1/2$, while all trajectories are holomorphic and
bounded by two on $|U|,|V|\le2$. If $b_{kl}$ are Taylor
coefficients of any such state or head, Cauchy's estimate gives
\[
 |b_{kl}|\le2^{1-k-l}.
\]
Retain the rectangle $0\le k,l\le J$. Its value error and
its derivatives of total order at most four, with respect to
$U,V$, are bounded by $C(J+2)^4 4^{-J}$. Conversion to the
original parameters costs at most $(2a)^{-4}$. Thus degree
$J=O_s(P+\log(r+2))$ and coefficient precision
$S=P+4\lceil\log_2(1/a)\rceil+O(\log(J+2))$ suffice for
absolute error $2^{-P}$ in the value and the derivatives
$f_{11},f_{22},f_{1122}$ needed by the two-variable correction.
Changing from means to Bernoulli biases changes these estimates
only by fixed factors.

We justify preprocessing these coefficient arrays without an
exponential-in-depth precision allowance. For polynomials
truncated to the rectangle, use the sum of absolute coefficients
as norm. Truncation preserves submultiplicativity. The exact
state jets have norm at most eight. All local analytic functions
and their required inverses have bounded coefficient norm
$O_s(1)$, by the same radius-two Cauchy estimate. In particular,
the squared denominator, its inverse, and its inverse square
root have such bounds; the argument remains in the right
half-plane bounded away from zero. The exponential and its
denominator in tanh also have bounded norms depending on $s$.
Neumann series, or the convergent analytic power series about
the exact jet, consequently show that each local operation is
$O_s(1)$-Lipschitz in this coefficient norm whenever the input
error is sufficiently small. For inverse square roots, this
follows explicitly by writing
$(V+E)^{-1/2}=V^{-1/2}(1+V^{-1}E)^{-1/2}$.

For explicit error budgets, one may take
\[
 K_s=2^{100(\lceil s\rceil+2)},\qquad C_s=K_s^{100},
 \qquad d_{\mathrm{err}}=12.
\]
Indeed, on the radius-two parameter polydisk the squared
denominator has modulus at most ten and real part at least
$1/24$. Its square root, inverse, and inverse square root
have modulus at most $24$. The argument of tanh has modulus
at most $10s$ and imaginary part at most $1/4$. Thus its
exponentials, their reciprocals, and the inverse denominator
of tanh have modulus bounded by $2e^{20s}+24$.
Cauchy's coefficient bound multiplies these supremum bounds
by at most four. The displayed $K_s$ bounds all these
coefficient norms; the elementary product and inverse
identities then give the stated larger $C_s$ as a common
local Lipschitz and residual-error constant.

At $M_*$ working bits, the polynomial procedures below
contribute at most $C_s(J+1)^{12}2^{-M_*}$ error.
The normalized update therefore
satisfies
\[
 E_{j+1}\le(1+C_sq_j)E_j+
                     C_s(J+1)^{12}2^{-M_*}.
\]
Its product is at most $r^{C_s}$, because $\sum q_j\le\log r$.
It follows that
\[
 E_D\le C_s(D+1)r^{C_s}(J+1)^{12}2^{-M_*}.
\]
For example, take an integer $M_*$ at least
\[
 S+\lceil\log_2(C_s(D+1))\rceil
   +\lceil C_s\log_2r\rceil
   +12\lceil\log_2(J+1)\rceil
   +\lceil\log_2(64K_s)\rceil+16.
\]
Certified dyadic upper enclosures of these logarithms determine
such an integer without testing transcendental equality.
This is $S+O_s(\log(D+2)+\log(r+2)+\log(J+2))$.
It ensures all promised final errors and, by induction, keeps the
intermediate errors inside the neighborhoods just used.
The known constant $c=\tanh1$ is enclosed to the same accuracy.
Stable evaluation of the known scalar coefficient $b_j$ uses
\[
 b_j=\eta c^2-\frac{2\eta c(c+w)}{r_j}
       +\frac{4-\eta z^2+2\eta cw+\eta c^2}{r_j^2},
\]
so its coefficient magnitudes remain bounded.

For completeness, the polynomial computations take
$O(J^4)$ coefficient operations per layer. Multiplication and
inversion use the elementary coefficient recurrences.
Let $\mathcal E=U\partial_U+V\partial_V$ be the Euler
operator, which multiplies a coefficient by its total degree.
It satisfies the product rule in the rectangular quotient
algebra and its inverse on zero-constant polynomials has norm
at most one. For $Y=\exp P$, solve
$\mathcal EY=Y\mathcal EP$ with $Y_{00}=e^{P_{00}}$;
explicitly, for nonzero multi-index $\lambda$,
\[
 |\lambda|Y_\lambda
 =\sum_{0<\mu\le\lambda}|\mu|P_\mu Y_{\lambda-\mu}.
\]
For $Y=V^{-1/2}$, compute $V^{-1}$ and
$H=-V^{-1}\mathcal EV/2$, then solve
$\mathcal EY=YH$ with $Y_{00}=V_{00}^{-1/2}$.
Here $H_{00}=0$. Both recurrences divide by total degrees
between one and $2J$, and use $O(J^4)$ operations.

These recurrences have polynomial residual error at fixed
precision. For the exponential, a residual $R$ obeys
$\mathcal E(e^{-P}Y)=e^{-P}R$; for the inverse square root
it obeys $\mathcal E(V^{1/2}Y)=V^{1/2}R$.
Apply $\mathcal E^{-1}$ after removing the constant term,
then multiply by the exact solution. The coefficient norms
bounded above show that the output error is at most a fixed
$C_s$ times the residual norm and the constant-term error.
Each residual coefficient sums $O(J^2)$ rounded products,
of magnitude at most $O_s(J)$, and there are $O(J^2)$
coefficients. The exponent $d_{\mathrm{err}}=12$ therefore dominates all
residual accumulation, including the fixed number of
intermediate multiplications and inversions. Inversion itself
uses the identity $VY-1=R$ and the bound on $V^{-1}$.
A simultaneous induction over total coefficient degree keeps
computed coefficient norms within one of their stated bounds:
apply the same residual identity in the further quotient that
discards all larger total degrees, and use the precision margin
above to make the error smaller than $1/(64K_s)$. This closes
the bounded-intermediate assumption used in charging products.
Direct fixed-precision arithmetic and certified scalar
exponential and square-root routines therefore take
$O_s((D+1)B^6)$ bit work per patch when $P=O(B)$.
The Taylor degree, coefficient precision, and all guard lengths
are $O_s(B)$. Nested evaluation of the stored polynomial and
its required derivatives at rational arguments uses
$O_s((B+\operatorname{bits}(z,w))^4)$ bit work.

Compute a rational $\Lambda$ with
$r^{2\nu}\le\Lambda\le2r^{2\nu}$ by certified logarithm
and exponential arithmetic. If $\Lambda\le n$, use the Taylor
table; otherwise use the finite lattice table described below.
This decision involves only rational comparisons.

In the first case $r^{2\nu}\le n$. Precompute both head
patch tables at precision
\[
 \epsilon=2^{-P},\qquad
 P=32+2\lceil\log_2(D+2)\rceil.
\]
This costs $O_s((D+1)nB^6)$. A corrected Bernstein coefficient
is a fixed linear combination of the value and the three
derivatives above, followed by at most one hypergeometric
average. The polynomial derivative prefactors have bounded
magnitude and the averaging weights are nonnegative and sum
to one. Hence cached evaluations determine the entire
coefficient to error $C\epsilon$, independently of its scale.
Rounding the rational prefactors and the averaging computation
to the same final tolerance uses the displayed polynomial bit
budget. For the hypergeometric average, use the mode-centered
recurrence in the proof of Theorem~\ref{thm:implicitrankone}:
set the mode weight to one and generate adjacent weights by
rational ratios. All weights are at most one, their sum is
at least one, and $O(\log(m+2))$ guard bits suffice.
This performs $O(m)$ arithmetic steps without constructing
binomial integers of length proportional to $m$.

Let $m_0=O_s(r^{2\nu})$ be the factory base size and let
$c_m=\overline A/(4m^2)$ be a correction mass, on dyadic
levels $m=2^km_0$. The known constant
$\overline A=O_s(m_0^2)$ is the sum of the derivative budgets
from the two-variable factory. Use the cached enclosure first
when comparing the normalized coefficient with its fresh
uniform random number. Conditional on selecting this
correction, a direct recurrence evaluation is necessary only
with probability at most $C\min\{1,\epsilon/c_m\}$.
If $c_m<\epsilon$, simply use the direct method immediately.
Thus its unconditional probability is at most
$C\min\{c_m,\epsilon\}$. We charge its refinements by
unconditional interval measures. If $c_m\ge\epsilon$, after
$t$ additional bits the unresolved set of uniform values has
measure at most $C(\epsilon/c_m)2^{-t}$. Computing that
enclosure costs
$O_s((D+1)m(B+\log(m+2)+P+t)^4)$.
Multiply by the branch mass $c_m$ and sum over $t\ge0$.
The resulting unconditional contribution is
$O_s((D+1)m\epsilon(B+\log(m+2)+P)^4)$.
If $c_m<\epsilon$, start the direct computation at absolute
precision proportional to $c_m$; the same argument gives
$c_m$ in place of $\epsilon$. This argument does not assume
that the work and the cache-miss event are independent.
Since $P=O(B)$, splitting the geometric sum at
$m\asymp\sqrt{\overline A/\epsilon}$ gives
\[
 \sum_{m=2^km_0}\min\{\overline A/m^2,\epsilon\}
             m(B+\log(m+2))^4
 =O_s\bigl(m_0\sqrt\epsilon\,B^4\bigr).
\]
Here $\log m_0$ and $\log(1/\epsilon)$ are $O_s(B)$; the
remaining polynomial factors have finite geometric sums.
Since $(D+1)\sqrt\epsilon\le1$, all direct fallbacks cost
$O_s(m_0B^4)$ in expectation. The base comparison costs less:
its fallback probability is $O(\epsilon)$.
The cached work and input probes themselves have expected
cost $O_s(m_0B^4)$ by the source-cost accounting of the factory.
Every comparison refines certified enclosures until it decides,
so this modification preserves the exact output law.

In the second case $r^{2\nu}>n/2$. Both unknown means belong to
the known lattice $\{-1,-1+2/N,\ldots,1\}$. Precompute the
two head values at all $(N+1)^2$ pairs, each to error
$\epsilon$. Direct scalar recurrence evaluation costs
$O_s((D+1)n^2B^4)$ for this table. Read the $2N$ relevant
input signs, retrieve the appropriate enclosure, and compare
with a fresh uniform. With probability $O(\epsilon)$ refine
by direct recurrence. Its expected extra cost is
$O_s((D+1)\epsilon B^4)=O_s(B^4)$, and the total online
cost is $O_s(nB^4)=O_s(\min\{r^{2\nu},n\}B^4)$.
Both preprocessing alternatives meet the
claimed budget. Upward rational enclosures decide all mesh
sizes and derivative budgets, as in the preceding theorem.
\end{proof}
